% !PRINTED 429
\weberpage{429}
%<<CONT
there follows from this the equation of degree 6 for $U$:
\weq{7}{U^6 - 10\,z\,U^3 + z^2\,U + 5\,z^2 = 0}
as the \gesperrt{resolvent of degree} 6 of the icosahedral equation.

If, following (3) and (4), we set
\weq{8}{\begin{aligned}
\sqrt{w_\infty} &= \sqrt{5}\,x_1\,x_2 \\
\sqrt{w_r} &= \varepsilon^r x_1^2 + x_1 x_2 - \varepsilon^{-r} x_2^2,
\end{aligned}}
then from this the following relations result:
\weq{9}{\begin{aligned}
\sqrt{5\,w_\infty} &= \sqrt{w_0} + \sqrt{w_1} + \sqrt{w_2} + \sqrt{w_3} + \sqrt{w_4} \\
0 &= \sqrt{w_0} + \varepsilon^2\sqrt{w_1} + \varepsilon^4\sqrt{w_2} + \varepsilon\sqrt{w_3} + \varepsilon^3\sqrt{w_4} \\
0 &= \sqrt{w_0} + \varepsilon^3\sqrt{w_1} + \varepsilon\sqrt{w_2} + \varepsilon^4\sqrt{w_3} + \varepsilon^2\sqrt{w_4}
\end{aligned}}
\weq{10}{\frac{x_1\sqrt{5}}{x_2} = \frac{\sqrt{w_0} + \varepsilon^4\sqrt{w_1} + \varepsilon^3\sqrt{w_2} + \varepsilon^2\sqrt{w_3} + \varepsilon\sqrt{w_4}}{\sqrt{w_\infty}},}
and in the equations (9), (10), since they are homogeneous, the $U$ may also be put in place of the $w$.

The equations of degree 6 whose roots satisfy the relations (9) were first pointed out by \gesperrt{Jacobi}; they are therefore called \gesperrt{Jacobi} equations of degree 6\footnote{Suite des notices sur les fonctions elliptiques. \gesperrt{Crelle's} Journal, Bd. III (1828); Werke, Bd. I, S. 261.}.

\begin{center}
\S.\,114.\\[4pt]
\gesperrt{On the solution of the icosahedral equation by transcendental functions.}
\end{center}

We wish here, without entering into proofs or more detailed developments, to point briefly to the relations that exist between the investigations we have carried out and the theory of certain transcendental functions, in particular the theory of elliptic functions and the theory of the hypergeometric series, from which results the solution of the icosahedral equation, and with it of all its resolvents, by transcendental functions. This leads out of the region of algebra, but stands to it in about the same relation as when one employs logarithms in the computation of root quantities, or the division of the angle in the case of the cubic equations.
